\documentclass[a4paper,11pt]{article}
\usepackage[a4paper, margin=2.5cm]{geometry}
\usepackage{listings}
\usepackage{parskip}
\usepackage{microtype}
\usepackage{hyperref}
\usepackage{xcolor}

\lstset{
	basicstyle=\ttfamily\small,
	breaklines=true,
	frame=single,
	backgroundcolor=\color{gray!8},
	commentstyle=\color{gray},
	morecomment=[l]{//},
}

\title{The data stored in \texttt{results} directory}
\date{}

\begin{document}
	
	\maketitle
	
	For a given integer $n$, two files are stored.
	\begin{itemize}
		\item \texttt{Grps-n} is stored in \texttt{results/transgrps}. It contains a
		complete list of the transitive subgroups of $\mathrm{Hol}(N)$ for all groups $N$
		of order $n$.
		
		\item \texttt{Results-n} is stored in \texttt{results}. It describes the equivalence
		classes of these transitive subgroups and records the number of corresponding
		structures.
	\end{itemize}
	
	The subgroups are stored as a sequence, say $X$.
	
	Each equivalence class $C$ is recorded as a tuple having the following shape:
	
\begin{lstlisting}
< a,  bool,
[
// may be multiple entries of this type for different values of m
[ <m, [ b ], [ c ], [ d ], e, f, g, h, i, j>
...
] 
>;
\end{lstlisting}
	
	The contents of $C$ are the following:
	
	\begin{itemize}
		\item $a$ = size of $C$.
		
		\item If \texttt{bool} is \texttt{true}, then:
		\begin{itemize}
			\item the skew bracoids are essentially skew braces, otherwise not;
			\item the corresponding field extensions are Galois, otherwise they are
			not normal.
		\end{itemize}
	\end{itemize}
	
	Let $N = \mathtt{SmallGroup}(n, m)$ where 
	$m \in [1\,..\,\mathtt{NumberOfSmallGroups}(n)]$.
	If transitive subgroups of $\mathrm{Hol}(N)$ contribute to $C$, 
	then the sequence $C[3]$ has an entry 
	$$\langle m,\,[b],\,[c],\,[d],\,e,\,f,\,g,\,h,\,i,\,j\rangle.$$
	Its contents are the following.
	\begin{itemize}
		\item The sequence $[b]$ records the index positions in the sequence $X$ of these
		transitive subgroups of $\mathrm{Hol}(N)$.
		
		\item The sequence $[c]$ records the index positions in $X$ of those transitive subgroups which
		correspond to Hopf--Galois structures giving a bijective Hopf--Galois
		correspondence.
		
		\item The sequence $[d]$ records the index positions in $X$ of 
those transitive 
		subgroups which correspond to almost classical structures.
		
		\item $e$ = the number of subgroups (counting conjugacy) contributed 
		by $N$. (If we sum the values of $e$ for all $N$ contributing 
		to $C$, then we obtain $a$, the size of $C$.)
		
		\item $f$ = the number of Hopf-Galois structures corresponding to $N$.
		\item $g$ = the number of these Hopf-Galois structures which 
		admit a bijective correspondence.
		\item $h$ = the number of these Hopf-Galois structures  which are almost classically Galois.
		\item $i$ = the number of skew bracoids corresponding to $N$.
		\item $j$ = the number of these skew bracoids which are almost classical.
	\end{itemize}
	
	We record one illustrative example.
	
	\begin{lstlisting}
> X := eval Read ("results/TransGrps/Grps-8");

> // Holomorphs of all five groups of order 8 contain 148 conjugacy classes of transitive subgroups

> #X;

148

> Y := eval Read ("results/Results-8");

> // there are 35 equivalence classes

> #Y;

35

> // one equivalence class

> C := Y[29];

> C;

<8, false, [
<3, [ 73, 76 ], [ 73, 76 ], [ 73 ], 2, 2, 2, 1, 2, 1>,
<4, [ 102, 104 ], [ 102, 104 ], [ 102 ], 6, 2, 2, 1, 2, 1>
]>
\end{lstlisting}
	From this data, we deduce the following.
	
	\begin{itemize}
		\item 
		The size of the equivalence class is 8.
		
		\item 
		The skew bracoids are not essentially skew braces and 
		the corresponding field extensions are not Galois. 
		
		\item 
		The two groups contributing to $C$ are {\tt SmallGroup}$(8, m)$ for $m= 3, 4$. Let these be denoted by $N_1$ and $N_2$ respectively.
		
		\item 
		Two conjugacy classes of transitive subgroups of $\mathrm{Hol}(N_1)$ 
		contribute to $C$; representatives are $X[73]$ and $X[76]$; the 
		number of subgroups in these classes is 2. 
		
		\item 
		Two conjugacy classes of transitive subgroups of $\mathrm{Hol}(N_2)$
		contribute to $C$; representatives  are $X[102]$ and $X[104]$; 
		the number of subgroups in these classes is 6. 
		
		\item Let $L/K$ be a separable extension associated to $C$. Then $L/K$ admits 2 Hopf--Galois structures of type $N_1$, corresponding to the transitive subgroups $X[73]$ and $X[76]$, and 2 Hopf--Galois structures of type $N_2$, corresponding to the transitive subgroups $X[102]$ and $X[104]$. All such Hopf--Galois structures admit a bijective correspondence, while only the structures associated to $X[73]$ and $X[104]$ are almost classically Galois.
		
		\item For $1 \leq i \leq 4$, there are actions $\odot_i$ such that the triples $(X[73],N_1,\odot_1)$, $(X[76],N_1,\odot_2)$, $(X[102],N_2,\odot_3)$ and $(X[104],N_2,\odot_4)$ each represent equivalence classes of skew bracoids (in particular, two classes have type $N_1$ and two classes have type $N_2$). The bracoids $(X[73],N_1,\odot_1)$ and $(X[102],N_2,\odot_3)$ are almost classical. Note that the stabilisers of the actions are pairwise isomorphic.
		
	\end{itemize}
	
\end{document}
